\documentclass{article}
\usepackage[UTF8]{ctex}
\usepackage{amsmath,amssymb,amsfonts}
\usepackage{enumitem}
\usepackage{geometry}
\geometry{a4paper, margin=1in}
% === 错题本排版优化（不参与题目标识） ===
\setlength{\parindent}{0pt}
\setlength{\parskip}{0.35em plus 0.10em minus 0.08em}
\setlength{\emergencystretch}{3em}
\setlist{itemsep=0.25em, topsep=0.40em, parsep=0pt}
\allowdisplaybreaks[2]
% === 错题本排版优化结束 ===
\begin{document}
\section*{首都师范大学2026年高等代数真题}
\section*{30．首都师范大学2026年高等代数试卷}
\begin{itemize}

\subsection{题目1}
  \item[1．] 设 \(n\) 为大于1的正整数，计算下列 \(n\) 级行列式
\[
\left|\begin{array}{cccc}
1+x_{1} & x_{1} & \cdots & x_{1} \\
x_{2} & 2+x_{2} & \cdots & x_{2} \\
\vdots & \vdots & & \vdots \\
x_{n} & x_{n} & \cdots & n+x_{n}
\end{array}\right|
\]

\subsection{题目2}
  \item[2．] 设 \(A, B, C\) 及 \(D\) 均为数域 \(\mathbb{P}\) 上的 \(n\) 级矩阵且 \(D\) 可逆，令
\[
M=\left(\begin{array}{cc}
O & A D \\
C & B
\end{array}\right) .
\]
用 \(r(A)\) 表示矩阵 \(A\) 之秩，证明：\(r(M) \geq r(A)+r(C)\) ．
\end{itemize}

\subsection{题目3}
3．取实数域 \(\mathbb{R}\) 上的 3 级矩阵

\[
A=\left(\begin{array}{ccc}
3 & -1 & 1 \\
-1 & 3 & 1 \\
1 & 1 & 3
\end{array}\right)
\]

\begin{itemize}
  \item[（1）] 求正交矩阵 \(C\) ，使得 \(C^{\prime} A C\) 为对角矩阵，其中 \(C^{\prime}\) 表示 \(C\) 的转置．
  \item[（2）] 求正定矩阵 \(B\) ，使得 \(B^{2}=A\) 。
\end{itemize}

\subsection{题目4}
4．设 \(A\) 为数域 \(\mathbb{P}\) 上的 \(n\) 级矩阵，\(\beta\) 是 \(n\) 维非零向量，假设 \(\alpha_{1}, \alpha_{2}, \cdots, \alpha_{s}\) 是齐次线性方程组 \(A X=0\)的一组基础解系，\(\alpha\) 是非齐次线性方程组 \(A X=\beta\) 的一个解。证明：向量组 \(\alpha_{1}, \alpha_{2}, \cdots, \alpha_{s}, \alpha\) 是线性无关的。

\subsection{题目5}
5．设 \(A\) 是 \(n\) 级半正定矩阵，\(D\) 是 \(n\) 级实矩阵，\(A^{*}\) 表示 \(A\) 的伴随矩阵。

\begin{itemize}
  \item[（1）] 证明：若 \(A\) 正定，则 \(D^{\prime} A^{*} D\) 半正定。
  \item[（2）] 当 \(A\) 非正定时，请问 \(D^{\prime} A^{*} D\) 是半正定吗？说明理由。
\end{itemize}

\subsection{题目6}
6．设 \(f(x), P(x) \in \mathbb{Q}[x]\) ，且 \(P(x)\) 是有理数域 \(\mathbb{Q}\) 上的不可约多项式，设 \(P(x)\) 有一个复数根 \(\alpha\) ，如果 \(\alpha\) 也是 \(f(x)\) 的复根，证明：在 \(\mathbb{Q}[x]\) 中，\(P(x) \mid f(x)\) 。

\subsection{题目7}
7．设 \(\mathcal{A}\) 是数域 \(\mathbb{P}\) 上的 \(n\) 维线性空间 \(V\) 上的一个线性变换，如果 \(\mathcal{A}\) 的矩阵可以对角化，证明：对 \(\mathcal{A}\)的任意不变子空间 \(W, \mathcal{A}\) 限制在 \(W\) 上的变换 \(\mathcal{A} \mid W\) 的矩阵也可以对角化。

\subsection{题目8}
8．设 \(f\) 是复数域上 \(n\) 维线性空间 \(V\) 上的线性变换，这里 \(n \geq 2\) ，它在一组基下的矩阵 \(A\) 为对角线元素为 1 的若尔当块

\[
\left(\begin{array}{ccccccc}
1 & 1 & 0 & 0 & \cdots & 0 & 0 \\
0 & 1 & 1 & 0 & \cdots & 0 & 0 \\
0 & 0 & 1 & 1 & \cdots & 0 & 0 \\
0 & 0 & 0 & 1 & \cdots & 0 & 0 \\
\vdots & \vdots & \vdots & \vdots & & \vdots & \vdots \\
0 & 0 & 0 & 0 & \cdots & 1 & 1 \\
0 & 0 & 0 & 0 & \cdots & 0 & 1
\end{array}\right)
\]

（1）求 \(f\) 的特征多项式，特征值及相应的特征子空间的维数。请回答 \(A\) 可以对角化吗？说明理由。\\
（2）证明：矩阵 \(A\) 与 \(A^{2}\) 相似．

\subsection{题目9}
9．设 \(V, W\) 是数域 \(\mathbb{P}\) 上的两个线性空间，其维数分别为 \(\operatorname{dim} V=n, \operatorname{dim} W=m, \sigma\) 是 \(V\) 到 \(W\) 的一个满射，且满足线性性，即

\[
\sigma(\alpha+\beta)=\sigma(\alpha)+\sigma(\beta), \sigma(k \alpha)=k \sigma(\alpha), \forall \alpha, \beta \in V, k \in \mathbb{P} .
\]

记 \(U=\{\alpha \in V \mid \sigma(\alpha)=0\}\) ，证明：\(U\) 是 \(V\) 的子空间，且 \(\operatorname{dim} U=n-m\) ．

\subsection{题目10}
10．设 \(V\) 是实数域 \(\mathbb{R}\) 上的一个 \(n\) 维线性空间，在 \(V\) 上定义一个二元实值函数，记为 \([\alpha, \beta]\) ，且满足：对任意的 \(\alpha, \beta, \gamma \in V, k \in \mathbb{R}\) ，有

\begin{itemize}
  \item[（i）] \([k \alpha, \beta]=k[\alpha, \beta]\) ．
  \item[（ii）] \([\alpha+\beta, \gamma]=[\alpha, \gamma]+[\beta, \gamma]\) ．
  \item[（iii）] \([\alpha, \beta]=-[\beta, \alpha]\) ．
  \item[（iv）] 如果 \([\alpha, \beta]=0\) 对任意的 \(\beta \in V\) 成立，则有 \(\alpha=0\) 。
\end{itemize}

此时我们称 \(V\) 关于该二元函数构成一个实数域 \(\mathbb{R}\) 上的 \(S\) 空间。证明：对于一个 \(n\) 为 \(S\) 空间 \(V\) ，以下结论成立：

\begin{itemize}
  \item[（1）] \(n \neq 1\) 。
  \item[（2）] 对于 \(V\) 中的任意非零向量 \(\alpha\) ，存在 \(\beta \in V\) ，使得 \([\alpha, \beta]=1\) 。
  \item[（3）] 设 \(K\) 为 \(V\) 中的由（2）中的 \(\alpha, \beta\) 生成的子空间，记
\[
K^{\perp}=\{\gamma \in V \mid[\gamma, \delta]=0, \forall \delta \in K\} .
\]
证明：\(K^{\perp}\) 也是 \(V\) 的子空间，且 \(V=K \oplus K^{\perp}\) 。
  \item[（4）] 证明：\(n\) 为偶数（记 \(n=2 k\) ），且存在 \(V\) 的一组基 \(\varepsilon_{1}, \varepsilon_{2}, \cdots, \varepsilon_{k}, \varepsilon_{-1}, \varepsilon_{-2}, \cdots, \varepsilon_{-k}\) ，使得
\[
\left[\varepsilon_{i}, \varepsilon_{-i}\right]=1, \forall 1 \leq i \leq k ;\left[\varepsilon_{i}, \varepsilon_{j}\right]=0, \forall i, j \in\{ \pm 1, \pm 2, \cdots, \pm k\} \text {, 并且 } i+j \neq 0 \text {. }
\]
\end{itemize}
\end{document}
